%% ---------------------------------------------------------------------
%% 04 - THẢO LUẬN
%% Diễn giải kết quả, so sánh với công trình trước, hạn chế của nghiên cứu.
%% ---------------------------------------------------------------------

\section{Thảo luận}

Kết quả cho thấy phần lớn thời gian tính toán nằm ở bước lấy mẫu chứ không
phải ở bước dựng biểu diễn hình học. Điều này phù hợp với quan sát trong
\cite{karnik2021} rằng chi phí của RRT* chủ yếu đến từ số lần kiểm tra va
chạm. Khi miền lấy mẫu đã được giới hạn bởi hành lang Voronoi, số lần kiểm
tra giảm rõ rệt mà không làm mất tính đầy đủ của thuật toán.

\emph{(Phần này là văn bản mẫu. Thay bằng nội dung thật: so sánh định lượng
với các phương pháp cơ sở, giải thích nguyên nhân của khác biệt, và nêu rõ
các hạn chế — ví dụ môi trường chỉ có vật cản lồi, chưa thử với nhiều robot,
chi phí tính toán trên máy nhúng.)}
